\centerline{\bf \Huge Начало числовых ребусов}

\begin{flushright}
        \scriptsize{\textbf{Разве ты не заметил, что способный к математике изощрен во всех науках в природе?\\
         Платон}}
        
\end{flushright}

\textbf{В ребусах с участием букв действует следующее правило: }

\textit{Одинаковым буквам должны соответствовать одинаковые цифры.}

\textit{Разным буквам — разные цифры.}


\textbf{Решить ребус - значит найти ВСЕ решения и доказать, что других нет!}\\

\problem{\textbf{1}} Ы + Ы + Ы + Ы + Ы + Ы + Ы + Ы + Ы = 90\\

\problem{\textbf{2}} 4А+5Б=100\\

\problem{\textbf{3}} А + Б + Б = 13\\

\problem{\textbf{4}} ОО + ОК = 67\\

\problem{\textbf{5}} СЕМЬ + СЕМЬ = 7777\\

\problem{\textbf{6}} НЕТ - НЕ  - Н = ДА\\

\problem{\textbf{7}} ТУТ + ТАМ = 222\\

\problem{\textbf{8}} АБВ + АВБ + ААА = 1000\\

\problem{\textbf{9}} КТО + КТО + КТО = ЧТД\\

\problem{\textbf{10}} ГОЛ + ЛОГ = ЧТО\\


